\section{A Formal Framework for Verbal Reinforcement Learning}
\label{sec:framework}

The methods surveyed in later sections share a decision-theoretic skeleton but attach language to it in different places. This section fixes that skeleton once. We define a language-augmented partially observable Markov decision process (POMDP), enumerate the points at which a linguistic artifact, written $\ell$ throughout, can enter it, and introduce an improvement operator $\mathcal{I}$ whose three realizations underpin the taxonomy of \S\ref{sec:taxonomy}. We then restate the two criteria that fence verbal reinforcement learning (VRL) off from adjacent paradigms, trace the genealogy of the term, and close by stating plainly which formal guarantees exist and which do not.

\subsection{The Language-Augmented Decision Process}
\label{subsec:setting}

We begin from a standard POMDP $\mathcal{M} = (\mathcal{S}, \mathcal{A}, \Omega, T, R, O, \gamma)$ \citep{kaelbling1998planning}: a Markov decision process with state space $\mathcal{S}$, action space $\mathcal{A}$, transition kernel $T$, reward function $R$, and discount $\gamma$ \citep{sutton2018reinforcement}, made partially observable by an observation space $\Omega$ and an observation function $O$ through which the agent perceives state only indirectly. Let $\mathcal{L}$ denote the set of finite natural-language strings. A \emph{linguistic artifact} $\ell \in \mathcal{L}$ is any string that the system produces or consumes as a signal: a task description, a critique, a rationale, a stored lesson, or a generated fragment of reward code. The agent is a configuration $\kappa = (\theta, c, m)$ comprising model parameters $\theta$, a working context $c$ (the prompt and any in-context material), and a persistent memory $m$; the policy $\pi_{\kappa}$ maps interaction histories to distributions over actions, which in language settings are themselves strings.

The construct that distinguishes VRL from ordinary language-conditioned control is the \emph{improvement operator}
\begin{equation}
\mathcal{I} : (\kappa, \ell) \longmapsto \kappa',
\label{eq:improvement}
\end{equation}
which consumes an artifact and returns a modified agent configuration. Three minimal realizations recur throughout the literature. An \emph{output revision} alters only $c$: a critique is appended to the context and the agent regenerates, as in iterative self-refinement \citep{madaan2023selfrefine}, and the effect vanishes when the episode ends. A \emph{memory write} alters only $m$: a reflection or distilled lesson is stored and retrieved in later episodes \citep{shinn2023reflexion,zhou2025memento}, so the effect persists without touching the weights. A \emph{parameter update} alters only $\theta$: the artifact, or a signal derived from it, changes the model itself and therefore all future behavior. Composite operators exist; Constitutional AI first revises outputs against written principles and then trains on the revisions, chaining the first and third realizations \citep{bai2022constitutional}.

Language can enter this augmented process at six points.
\begin{enumerate}[label=(\roman*),leftmargin=2.2em,itemsep=0.2em]
    \item \emph{Goal specification.} A task description $g \in \mathcal{L}$ defines what counts as success, often without any executable reward; the true objective is latent and reaches the agent only through language, the setting formalized as learning from language feedback in \S\ref{subsec:guarantees}.
    \item \emph{State and observation.} Observations may be linguistic (text environments, tool outputs, error traces) or verbalized descriptions of non-linguistic state; language then acts as a state abstraction whose fidelity bounds what the policy can distinguish (\S\ref{sec:grounding}).
    \item \emph{Action prior.} Instructions and subgoals reshape the distribution over actions before any learning occurs; in hierarchical designs an LLM planner emits subgoals in $\mathcal{L}$ that a low-level actor must realize \citep{he2024words}.
    \item \emph{Reward.} Timing is essential here. At \emph{specification time}, a generator compiles $\ell$ into an executable reward function $\hat{R} = C(\ell)$, as in reward-code generation \citep{ma2024eureka}: the linguistic structure is consumed by the compiler even though the resulting function emits scalars. At \emph{runtime}, a judge model maps trajectories to evaluations, verbal or scalar, on the fly \citep{zheng2023judging}. \S\ref{sec:rewardstack} develops this stack in full.
    \item \emph{Transition prior.} Pretrained linguistic knowledge supplies an implicit world model: predictions about the consequences of an action can be elicited in language before, or instead of, environment interaction \citep{he2024words}.
    \item \emph{The improvement operator.} Critiques, reflections, and preference rationales serve as the argument $\ell$ in Eq.~\eqref{eq:improvement}; this entry point hosts the second and third pillars of the taxonomy.
\end{enumerate}

The taxonomy of \S\ref{sec:taxonomy} classifies a signal by which of these points it acts on and, for the last, by which realization of $\mathcal{I}$ consumes it. Signals that instantiate the goal, observation, action-prior, or reward components of $\mathcal{M}$, points (i) through (iv), act at problem-definition time and constitute the grounding signal of \S\ref{sec:grounding}; artifacts consumed by output revision or memory write at inference time constitute deliberative feedback (\S\ref{sec:deliberative}); artifacts distilled into parameter updates constitute the learning signal (\S\ref{sec:learning}). Point (v) is the exception, and we state it plainly: a transition prior is not written down when the problem is defined but elicited from the pretrained model during interaction, and no method we survey consumes a linguistic artifact at this point. The taxonomy therefore operationalizes five of the six entry points; the transition prior enters the survey only through the theoretical analysis of \S\ref{subsec:guarantees} and stands as an entry point the literature has yet to occupy. Two caveats keep the formalism honest. First, the framework is an organizing lens, not a claim that every method admits a POMDP formulation, and the boundary is concrete. Open-ended web and tool-use tasks specify no transition kernel and no reward function; success is adjudicated post hoc by a judge or a human, so $T$ and $R$ are notational conveniences rather than defined objects. Dialogue sets the agent against an interlocutor whose state admits no compact Markovian summary. And once memory writes alter $m$, later episodes unfold in an environment the agent itself has changed, breaking stationarity (\S\ref{sec:credit}). In these regimes we use the formalism as vocabulary and do not import its theorems. Second, multi-turn language interaction typically carries a two-level structure, an utterance-level decision process layered over a token-level one, and both benchmark design and algorithm design increasingly make this decomposition explicit \citep{abdulhai2023lmrl,zhou2024archer}.

\subsection{The Definitional Fence}
\label{subsec:fence}

With this notation, the definition of \S\ref{sec:intro} can be restated precisely. A method belongs to VRL if and only if (i) the feedback is expressed in natural language, whether self-generated, provided by a human, or produced by an external tool or model, and its linguistic structure is functionally consumed by the agent rather than being collapsed to a scalar or preference bit before use; and (ii) the feedback participates in an improvement loop, that is, an output revision, a memory write, or a parameter update.

Criterion (i) can be phrased through a scalarization map $\sigma : \mathcal{L} \to \mathbb{R}^{k}$: if replacing $\ell$ by $\sigma(\ell)$ everywhere leaves the method's behavior unchanged, the linguistic structure was not functionally consumed and the criterion fails. Two terms in this test need pinning down. \emph{Unchanged} means distributional, not token-level, invariance: the distribution over updated configurations $\kappa'$, and hence over subsequent behavior on the method's own task distribution, is the same under $\ell$ and $\sigma(\ell)$. \emph{Small $k$} means that $k$ is fixed by the method rather than growing with the content of $\ell$, and it is counted per judged unit: a preference bit, a grade, a fixed vector of rubric scores, or one score per reasoning step is a scalarization, however many steps the judged trajectory contains; retaining the text is not. Criterion (ii) requires that $\ell$ appear as the argument of $\mathcal{I}$, or construct a component of $\mathcal{M}$ that subsequently drives $\mathcal{I}$, as a compiled reward function does.

The test is still a counterfactual, and \S\ref{subsec:guarantees} concedes that no measured quantity certifies it, so we state the inspection procedure used to adjudicate membership throughout this survey. For each method we trace the artifact from its producer to the operator that consumes it and mark the \emph{scalarization boundary}, the point at which language collapses into the numbers the learner receives; \S\ref{sec:rewardstack} develops this construct in full. What crosses the interface to $\mathcal{I}$ decides the case. Criterion (i) holds when some consumer downstream of the artifact reads its structure: a compiler that parses $\ell$ into reward code \citep{ma2024eureka}, a reviser that conditions on a critique \citep{madaan2023selfrefine}, a retriever that matches stored lessons \citep{shinn2023reflexion}, or a training loss computed over the text itself. It fails when the only consumer is a judge or annotator whose output is the scalar. The intermediate formats that dominate practice resolve by the same trace. A rationale-annotated preference qualifies when the rationale enters the training distribution or a revision loop, and does not when only the bit survives. A step-level critique qualifies when downstream machinery converts its content into credit, as when critique text is converted into token-level credits \citep{xie2025capo}, and reduces to process supervision when only per-step scores survive \citep{lightman2024let}. A generative process reward model reasons in language before scoring \citep{zhao2025genprm}; its boundary sits at verdict emission, so the linguistic structure is consumed inside the reward computation while the optimizer receives per-step scalars, the same route by which compiled rewards enter, unless the critique is separately routed into revision or critique-level supervision \citep{wang2026reward}, which moves the boundary later. The trace reduces the counterfactual to a question about dataflow; the residual judgment is empirical, namely whether a consumer positioned to read the structure in fact exploits it, and only the missing certificate of \S\ref{subsec:guarantees} would replace that judgment with a measurement. \S\ref{sec:learning} arranges learning-signal methods along the resulting compression spectrum, from methods that keep $\ell$ verbatim in the training distribution to methods that retain only $\sigma(\ell)$.

Both criteria are required, and each excludes a neighboring paradigm. Plain instruction following fails criterion (ii): the instruction conditions the policy, but nothing is revised, written, or updated, so no improvement loop exists. Vanilla RLHF \citep{ouyang2022training} and DPO \citep{rafailov2023direct} fail criterion (i): the annotator's verbal judgment is collapsed to a preference bit before it reaches the update, and the content that would make the signal verbal, the reason one response beats another, is discarded; we treat these methods only as the scalar endpoint of the compression spectrum. Language-conditioned RL \citep{luketina2019survey} treats language as an input to the policy rather than as feedback on its behavior; improvement is driven by an ordinary scalar reward, so language never enters the loop of criterion (ii). One boundary case the criteria settle explicitly: language consumed at specification time to construct a reward function satisfies criterion (i), because the generator exploits the full structure of $\ell$, and criterion (ii) through the training loop the constructed reward drives \citep{ma2024eureka}. Finally, the criteria fix the modality scope: the feedback signal is natural-language text, while observations may be non-textual, as in embodied settings where verbal corrections accompany visual state (\S\ref{sec:embodied}).

\subsection{Genealogy and Neighboring Formalisms}
\label{subsec:genealogy}

The phrase predates machine learning: in mid-twentieth-century behavioral psychology, verbal reinforcement denoted approving utterances an experimenter used to shape a subject's responses \citep{greenspoon1955reinforcing,krasner1958studies}. In the LLM era the term was introduced by \citet{shinn2023reflexion}, whose subtitle, ``Language Agents with Verbal Reinforcement Learning,'' named a specific mechanism: agents reflect verbally on task feedback and keep those reflections in an episodic buffer that guides later attempts, an instance of $\mathcal{I}$ realized as a memory write. We generalize the term deliberately to any method satisfying the two criteria above, flagging the departure from the original, narrower sense.

The closest formal ancestor of our framework is Natural Language Reinforcement Learning \citep{feng2024natural}, which extends the core objects of RL into language counterparts: the language value function replaces a scalar estimate with a narrative articulating the rationale behind an evaluation, and Bellman backups and policy iteration are redefined analogously, with LLMs implementing the resulting operators. Its authors argue that scalar value restricts an agent's understanding of its environment and report gains in effectiveness and interpretability on multi-step agentic tasks. In our notation, NLRL relocates the codomain of the value function from $\mathbb{R}$ to $\mathcal{L}$; the contribution is a framework, not an overview of the field, and the language-space optimizers of \S\ref{sec:optimizers} build directly on this move.

A second neighbor replaces reward maximization with conditional generation: the feedback-conditional policy \citep{luo2025fcp} learns the posterior over responses given feedback and bootstraps online under positive feedback conditions. Its motivating argument, that scalarizing nuanced feedback discards richness, is the probabilistic form of criterion (i): the feedback string survives into the training objective itself.

A third neighbor makes memory the learned object: Memento \citep{zhou2025memento} casts continual adaptation as a memory-augmented Markov decision process (M-MDP) in which the policy improves by rewriting and retrieving episodic cases, with the underlying LLM never fine-tuned. The M-MDP makes precise the sense in which memory is persistent state rather than inference-time scaffolding, the boundary question experiential memory poses for the taxonomy (\S\ref{sec:taxonomy}); the credit-assignment problem that memory writes create \citep{yan2026memoryr2} returns in \S\ref{sec:credit}.

Two earlier organizing efforts situate ours. Language grounding in RL predates LLMs \citep{branavan2009reinforcement}, and \citet{luketina2019survey} surveyed that first wave, dividing it into language-conditional RL, where language is part of the problem, and language-assisted RL, where language aids learning. Our framework is the LLM-era successor to that division: language is no longer only an input or an auxiliary knowledge source but the medium of evaluation and revision, and the policy itself is a language model, so every entry point of \S\ref{subsec:setting} is realizable by a single substrate. CoALA \citep{sumers2023cognitive} instead organizes agent computation into memory modules, action spaces, and decision cycles. The two lenses are complementary: CoALA says where computation and storage live inside an agent, whereas we classify the signals that flow through it, and a single CoALA-style architecture can host grounding, deliberative, and learning signals simultaneously (\S\ref{sec:taxonomy}).

\subsection{What Is Formally Guaranteed, and What Is Not}
\label{subsec:guarantees}

Formal results exist, but they are recent and localized. \citet{xu2025formalizing} give the first principled formalization of the learning-from-language-feedback (LLF) problem: sufficient assumptions under which learning is possible even though rewards are latent, a complexity measure, the transfer eluder dimension, that quantifies how informative feedback is, and a no-regret algorithm, HELiX, whose guarantees scale with that dimension. Their separation result is the theoretical warrant for the field's premise: learning from rich language feedback can be exponentially faster than learning from reward alone. Complementing this, \citet{he2024words} analyze a hierarchical planner and actor system inside a POMDP and prove that, under assumptions on the pretraining distribution, an LLM planner performs Bayesian aggregated imitation learning through in-context learning; naive execution of the imitated subgoals incurs linear regret, and epsilon-greedy exploration on top of the imitator restores sublinear regret when the pretraining error is small. The lesson generalizes: language-mediated hierarchies do not escape the exploration problem. At the level of problem statements rather than theorems, LLF-Bench operationalizes the setting, replacing numeric reward with natural-language instructions and feedback and randomizing verbalizations so that agents must genuinely learn from the feedback channel \citep{cheng2023llfbench}.

The list of what is missing is longer. There is no convergence or contraction analysis for language-space Bellman operators, and to our knowledge no analogue of the policy improvement theorem exists for $\mathcal{I}$ in any of its realizations: nothing certifies that a revision, a memory write, or a critique-driven update makes the policy better rather than differently biased. Existing regret results assume that feedback carries reliable information about a latent ground-truth objective; no bound covers self-generated feedback that shares the biases of the policy it evaluates, the regime in which most deliberative methods operate (\S\ref{sec:safety}). The separation of \citet{xu2025formalizing} contrasts the ends of the compression spectrum, rich feedback against scalar reward; no complexity characterization exists for the intermediate formats that dominate practice, step-level critiques or rationale-annotated preferences. Criterion (i) itself lacks a certificate: no measurable quantity attests how much of an artifact's linguistic structure a learner actually exploits, so the fence must be applied by the boundary trace of \S\ref{subsec:fence}, which settles what a mechanism could read but not what it in fact uses. And no guarantee yet addresses noisy, adversarial, or strategically generated feedback channels. These gaps recur as concrete agenda items in \S\ref{sec:practitioner}.

The vocabulary of this section, $\ell$ for the linguistic artifact, $\kappa$ for the agent configuration, $\mathcal{I}$ for the improvement operator, and $\sigma$ for scalarization, does most of its work here, in the entry points, the fence, and the guarantees above; the later chapters inherit the distinctions rather than the symbols. The taxonomy of \S\ref{sec:taxonomy} rests on the mapping of \S\ref{subsec:setting}: entry points (i) through (iv) become its four grounding cells, the three realizations of $\mathcal{I}$ separate deliberative feedback from the learning signal, and the transition prior stands as the unoccupied sixth coordinate. The notation returns where an argument needs it, as $\sigma(\ell)$ does in the fence adjudication of \S\ref{sec:learning} and the credit-assignment analysis of \S\ref{sec:credit}; elsewhere the pillar chapters carry these distinctions in prose.
